%!TEX root = ../Master_template.tex
\chapter{Introduction}

\section{Background and Motivation}

The classical Traveling Salesman Problem (TSP) is a well-known problem in many fields,
including computer science, operations research, and applied mathematics. The main goal
is to find the shortest path along which a salesman must visit each city once and return
to the depot. TSP belongs to the class of NP-hard problems. In the real world, most
domains that use TSP include logistics, supply chain management, manufacturing,
telecommunications, and many more industries.

The classical TSP assumes that travel costs are static and symmetric. In reality, travel
times between two locations can vary depending on the departure time due to factors such
as traffic congestion, rush hours, weather conditions, maintenance, and others. This
motivates the study of the Time-Dependent Traveling Salesman Problem (TDTSP), where the
travel time for each edge depends on the departure time from its origin node.

Figure~\ref{fig:tdtsp_example} illustrates this concept with a simple example. A
salesman at the depot must visit two customers, A and B, and return to the depot. The
three panels show how travel times between the same locations change throughout the
morning. At 08:00 during rush hour, the travel time from the depot to customer A is 60
minutes (red arrow) while the route to B takes 65 minutes (black arrow). By 09:00, the
travel time from A to B increases to 75 minutes due to shifting congestion patterns. At
10:00, as rush hour subsides, most travel times vary further. This temporal variation
has important implications for route planning. A visiting sequence that is optimal when
departing at 08:00 may be suboptimal when departing at 09:00 or 10:00. An effective
routing strategy must therefore anticipate how traffic conditions will evolve during the
tour, not merely react to current conditions.

\begin{figure}[H]
    \centering
    \includegraphics[width=\textwidth]{figs/TDTSP_Example}
    \caption{Illustration of time-dependent travel times in TDTSP. The salesman at the
    depot (van) must visit customers A and B. Travel times (in minutes) vary based on
    departure time: 08:00 (rush hour), 09:00 (mid-morning), and 10:00 (off-peak). Red
    arrows indicate one direction of travel; black arrows indicate the reverse direction.}
    \label{fig:tdtsp_example}
\end{figure}

One of the earliest studies on the TDTSP used integer programming to solve it~\citep{picard1978}.
This foundational work inspired subsequent research in this area. Later, \citet{malandraki1992} established a foundation for this problem class
by proposing formulations for time-dependent vehicle routing problems (TDVRP) with
step-function travel times. Most studies on the TDTSP assume the First-In-First-Out
(FIFO) property. This property ensures that leaving later from a node cannot result in
arriving earlier at the destination (no overtaking).

Heuristic methods such as time-expanded formulations, simulated annealing, ant colony
optimization, and genetic algorithms can produce high-quality results, but solving the
TDTSP is a computationally expensive task. This was demonstrated by \citet{cordeau2014},
who showed that even sophisticated branch-and-cut and branch-cut-and-price algorithms
can only solve moderate-size instances to optimality.

\section{Deep Reinforcement Learning for Combinatorial Optimization}

In the past decade, deep reinforcement learning (DRL) and attention-based neural
architectures have been used to solve routing problems. The Pointer Network is an
encoder-decoder model that uses Sequence-to-Sequence (Seq2Seq) models to solve
routing problems~\citep{vinyals2015pointer}. \citet{bello2017}
replaced the Pointer Network's supervised learning (which relies on labelled training
data) with the REINFORCE algorithm using a simple exponential moving average baseline.

The introduction of the Transformer architecture by \citet{vaswani2017} significantly
revolutionized the field, as it replaced recurrent neural networks (RNNs) entirely with
self-attention mechanisms. \citet{kool2019} combined all of these innovations and
produced the Attention Model (AM), which has become the standard baseline for solving
routing problems. By following the footsteps of the Attention Model, \citet{zhang2023}
proposed a deep reinforcement learning (DRL) framework for variants of the
time-dependent TSP. They proposed two model variants called M1 and M2 to handle the
time dependent dynamic TSP. In a fixed customer base problem, all customers are known
at the start of the tour and remain unchanged throughout the whole tour. In contrast, a
dynamic customer base allows new customers to appear or existing customers to be removed
during the tour. M1 is more suitable for a fixed customer base, while M2 is better
suited for a dynamic customer pool. In this thesis, we use the M1 model as the
baseline, as we solve the deterministic TDTSP with a fixed customer set~\citep{zhang2023}.
Furthermore, the paper reports that M1 achieves more than 5\% improvement over the M2
model for a fixed customer set. However, we identified several gaps in the M1 model
that leave room for potential improvements.

\section{Research Gaps and Contributions}

This section identifies three significant gaps in the current state-of-the-art model
M1, and for each gap, we describe the limitation, the importance of filling the gap,
and our proposed solution.

\subsection{Gap 1: Lack of Awareness of the Tour Progress in the Decoder}

At each decoding step, the standard M1 decoder receives only the following information:
(1) graph embedding, (2) depot embedding, (3) a binary visited mask indicating which
nodes have been visited, (4) the embedding of the last-visited node, and (5) the
current traffic conditions. Crucially, the decoder has no explicit signal indicating
how far through the tour it has progressed. This lack of progress awareness may be
problematic for the TDTSP. In TDTSP, decisions made at different stages of the route
have fundamentally different characteristics and should therefore be treated differently.
For example, a tour beginning in the morning must anticipate afternoon traffic patterns,
while an evening tour may benefit from declining congestion as it progresses. Without an
explicit step ratio or temporal progress signal, the decoder cannot differentiate early
from late tour decisions.

Previous work by \citet{wu2021} recognized this limitation and proposed augmenting the
decoder with a Gated Recurrent Unit (GRU). At each decoding step $t$, the GRU takes
the current context vector as input and combines it with its hidden state from the
previous step ($H_{t-1}$), producing an enhanced context vector $H_t$ that implicitly
encodes the history of visited nodes and decisions made so far. While the GRU is
effective, it can be computationally expensive, and it is therefore worthwhile to
replace it with a lightweight MLP (Step-MLP) with tour-aware features (step ratio, last
3 visited nodes, visited node embeddings, etc.). We also conduct experiments by varying
the position of the MLP layer:

\begin{itemize}
    \item \textbf{Pre-attention MLP:} Nonlinear refinement of the query before
    attention computation.
    \item \textbf{Post-attention MLP:} Refinement of the glimpse vector after attention
    but before final logit computation.
\end{itemize}

Additionally, we improve accuracy by scaling the attention logits with a
scaling factor (temperature scaling factor) learned by a lightweight MLP layer
(temp-MLP).

\subsection{Gap 2: Absence of Classical Heuristic Knowledge in Neural Routing}

Classical routing heuristics (nearest neighbor, greedy cost minimization) often provide
reasonable, interpretable, and computationally inexpensive solutions. Therefore, we
incorporate this knowledge into the decoder and produce a hybrid approach.

Recently, researchers have begun exploring this direction. As a result, \citet{ye2023}
proposed DeepACO, where neural networks generate heuristic measures to guide Ant Colony
Optimisation search. More recently, Huang et al.\ (2025) presented ReLD, which adds a
logarithmic distance term to the attention logits:

\begin{equation}
    p_i = \operatorname{Softmax}\!\left(C \cdot \tanh\!\left(
        \frac{q_c^\top H_t}{\sqrt{d_h}} - \log(\mathrm{dist}_i)
    \right)\right)_{\!i}
    \label{eq:reld}
\end{equation}

However, this approach has two main limitations for time-dependent routing problems.
First, it uses static distance metrics (Euclidean distance) that are not suitable for
TDTSP. Second, the heuristic integration is fixed (simple logarithmic transformation)
without learnable parameters to adapt the heuristic influence based on problem state.

To address these limitations, this thesis proposes a Cost-Aware Gating mechanism that:
(1) uses time-dependent travel costs rather than static distances as the heuristic
signal, (2) employs a learnable gating parameter $\lambda_{\text{heur}}$ that can adapt
the heuristic influence, and (3) explores nonlinear transformations (piecewise-linear,
exponential) of heuristic logits. We blend the heuristic logits as a bias into the
attention logits via:

\begin{equation}
    \ell_i^{\text{blended}} = \ell_i + \lambda_{\text{heur}} \cdot h_i(s_t)
    \label{eq:blended_logit}
\end{equation}

where $\ell_i$ denotes the raw attention logit for candidate node $i$, $h_i(s_t)$ is
the heuristic score, and $\lambda_{\text{heur}}$ is a learnable parameter.

\subsection{Gap 3: Static Encoder Consuming the Full Traffic History}

The encoders of standard attention models, including our M1 baseline \citep{zhang2023},
take the full-day traffic profile as input. However, in real-world dynamic routing
problems, specific opening and closing times mean that feeding the full-day traffic
profile wastes both computational resources and time. Therefore, we feed a time
window, instead of the full history, into the encoder. We also refresh the
encoder when the current decision time exceeds the encoded window boundary. We could
not find any prior work in this area, and this thesis addresses this gap with the Time
Slicing with Safe Refresh mechanism.

\section{Thesis Structure}

The remainder of this thesis is organised as follows:

\textbf{Chapter 2: Related Work} provides a comprehensive review of the literature on
classical and time-dependent TSP, traditional heuristic methods, neural combinatorial
optimization, and deep reinforcement learning for routing problems.

\textbf{Chapter 3: Methodology} discusses the technical details of our approach. First,
we describe the baseline M1 model's architecture, including the encoder-decoder
attention model. Next, we discuss the implementation details of the three proposed
enhancements (Step-MLP, the cost-aware gating mechanism, and time slicing with safe
refresh). Finally, we describe the training procedure with the REINFORCE algorithm with
rollout baseline.

\textbf{Chapter 4: Experiments} presents the end-to-end experimental setup and results.
We describe the real-world Beijing dataset, baseline comparisons, and evaluation
metrics. We then present results for each proposed enhancement, including ablation
studies and hyperparameter sensitivity analyses. We examine computational costs,
training stability, and scaling behaviour from 19-node to 49-node instances.

\textbf{Chapter 5: Conclusion} summarises the contributions of this thesis and their
implications for the broader field of neural combinatorial optimization.
